\documentclass[11pt]{article}
\usepackage[T1]{fontenc}
\usepackage{lmodern,amsmath,amssymb,amsthm,booktabs,geometry,hyperref}
\geometry{margin=1in}
\hypersetup{hidelinks}
% Non-printing provenance annotations; no Lean coverage is claimed.
\newcommand{\coverage}[1]{}
\newcommand{\lean}[1]{}
\newcommand{\uses}[1]{}
\newcommand{\proves}[1]{}
\newcommand{\imports}[1]{}
\newcommand{\evidence}[1]{}

\newtheorem{theorem}{Theorem}[section]
\newtheorem{proposition}[theorem]{Proposition}
\newtheorem{lemma}[theorem]{Lemma}
\newtheorem{corollary}[theorem]{Corollary}
\newcommand{\Fp}{\mathbb F_p}
\newcommand{\one}{\mathbf 1}
\newcommand{\wt}{\operatorname{wt}}
\newcommand{\rk}{\operatorname{rank}}
\newcommand{\Tr}{\operatorname{Tr}}
\newcommand{\im}{\operatorname{im}}
\newcommand{\GL}{\operatorname{GL}}
\newcommand{\PGL}{\operatorname{PGL}}
\newcommand{\PSL}{\operatorname{PSL}}
\title{Strength-two trades and transversal cubic gates:\\the Clebsch cubic-phase codes and their magic}
\author{}
\date{Draft, September 2026}
\begin{document}
\maketitle
\begin{abstract}
A signed point configuration whose moments vanish through degree two defines a
CSS code with a transversal cubic phase. We apply this dictionary to conic
matchings, obtaining error-detecting codes with invariant-theoretic logical
phases, and give an elementary $[[2p,p-1,2]]_p$ construction for every prime
$p\geq5$. The Pauli spectrum of a homogeneous cubic phase is determined by its
Hessian-rank distribution. Exact finite censuses distinguish the conic resources
from products of elementary cubic states; the ten-qudit example over $\mathbb F_{11}$
is not Clifford-equivalent to a product across any bipartition. For the six-qudit
example over $\mathbb F_7$, a native fourteen-input factory costs fewer than
$16.116$ raw resources per accepted block, compared with lower bounds of $54$
and $36$ for two explicitly specified separate-distillation menus at the same
target block infidelity. A five-dimensional restriction identifies a chordal
cubic and gives the two associated sign choices distinct operational meanings.
\end{abstract}
\section{Signed trades as quantum codes}\label{sec:trades}
The same cancellation that makes two point configurations agree on quadratics
can make a cubic phase insensitive to a stabilizer coordinate. The surviving
third moment is then a logical gate. Our examples use perfect matchings of a
conic, but the cancellation mechanism applies to any signed configuration.

Fix a prime $p\geq5$ and write $\omega=\exp(2\pi i/p)$. A qudit has basis
$\{|x\rangle:x\in\Fp\}$, with $X(a)|x\rangle=|x+a\rangle$ and
$Z(b)|x\rangle=\omega^{bx}|x\rangle$. A \emph{strength-two signed trade}
is a list of distinct points $x_i\in\Fp^k$, with signs $\epsilon_i\in\{1,-1\}$,
such that
\begin{equation}\label{eq:moments}
 \sum_i\epsilon_i=0,\qquad \sum_i\epsilon_i x_i=0,\qquad
 \sum_i\epsilon_i x_ix_i^T=0.
\end{equation}
These equalities are in $\Fp$. Let $E$ have rows $x_i^T$ and let
$L=\{a\one+Eu:a\in\Fp,\ u\in\Fp^k\}\subseteq\Fp^n$.
The coordinatewise product span is denoted by $L^{\circ2}$.

\begin{theorem}[Trade-to-code dictionary]\label{thm:trade-code}
\coverage{absent}
Suppose $n=2p$, $k=p-1$, and $\dim L=p$. Then the CSS code with
X-stabilizer space $S=\langle\one\rangle$ and Z-stabilizer space $L^\perp$
has parameters $[[2p,p-1,2]]_p$. The signed transversal operation
$\bigotimes_i M_{\epsilon_i}$, where $M_c|z\rangle=\omega^{cz^3}|z\rangle$,
induces the logical diagonal gate
\begin{equation}\label{eq:logical}
 U_F|u\rangle=\omega^{F(u)}|u\rangle,\qquad
 F(u)=\sum_i\epsilon_i(x_i^Tu)^3.
\end{equation}
More generally, for $S\subseteq L$, a physical cubic with weights $w_i$
is constant on every coset of $S$ in $L$ if and only if
$\sum_i w_i s_i b_i c_i=0$ for all $s\in S$ and $b,c\in L$.
For $S=\langle\one\rangle$ this condition is $w\perp L^{\circ2}$.
\end{theorem}
\begin{proof}
Put $D=\operatorname{diag}(\epsilon_i)$. Equation~\eqref{eq:moments}
says that $L$ is isotropic for the nondegenerate pairing $v^TDw$.
Its dimension is half the length, so $L^\perp=DL$. In particular
$S\subseteq L$, and the CSS dimension is $\dim L-\dim S=p-1$.
An encoded basis is
\[
 |u_L\rangle=p^{-1/2}\sum_{a\in\Fp}|a\one+Eu\rangle.
\]
Expanding $\sum_i\epsilon_i(a+x_i^Tu)^3$ removes all terms involving $a$
by~\eqref{eq:moments}, leaving~\eqref{eq:logical}.

For the distance, no weight-one Z error lies in $S^\perp$.
Since the points are distinct, $e_i-e_j\in S^\perp\setminus L^\perp$
for every $i\ne j$, so $d_Z=2$. A weight-one vector cannot lie in $L$:
its pairing with $\one\in L$ under $D$ would be nonzero. Thus $d_X\geq2$.

For the general criterion, write $T_w(s,b,c)=\sum_i w_i s_i b_i c_i$.
The difference of the physical cubic at $z+s$ and $z$ is
$3T_w(s,z,z)+3T_w(s,s,z)+T_w(s,s,s)$. Its quadratic part must vanish
as a polynomial function on $L$; degrees are less than $p$.
Polarization, using $2,3\ne0$, gives $T_w(s,L,L)=0$.
This condition also kills the lower terms because $s\in L$, proving sufficiency.
\proves{thm:trade-code}
\end{proof}

This is a specialization of the signed qudit orthogonality mechanism in
Watson--Campbell--Anwar--Browne~\cite[Definition 2 and Lemmas 6--8]{watson},
with related divisible-code and prime-qudit constructions
in~\cite{haah,cab,krishna}. The contribution considered here is the choice of
configuration and the resulting logical resource. Our positioning follows the
bounded comparison recorded in ledger rows N1--N4 of the accompanying literature
audit; no classification of all CSS codes with these parameters is asserted.
The Schur-square conductor also governs the MDS--CSS transversal-group
classification~\cite{mds}. The distinction is the full product
$S\circ L\circ L$ in Theorem~\ref{thm:trade-code}: Schur-square codimension
one alone does not imply a non-Clifford gate. Here the X-space is only
$\langle\one\rangle$.

\subsection{An elementary family at every prime}
An example separates parameter existence from the conic geometry.
Let $V=\Fp^p/\langle\one\rangle$, write $a_i$ for the image of the $i$th
standard basis vector, and set $t=a_0$. The functional
$\sigma([z])=\sum_i z_i$ is well-defined because $p=0$ in $\Fp$.
Take the positive sheet to be $\{a_i\}$ and the negative sheet to be
$\{a_i+t\}$.

\begin{proposition}[Translation trade]\label{prop:translation}
\coverage{absent}\uses{thm:trade-code}
These sheets define a strength-two trade with affine span $V$ and
$\dim L=p$. Its code is $[[2p,p-1,2]]_p$, and its logical cubic is
\begin{equation}\label{eq:translation-phase}
 F(u)=-3u(t)q(u),\qquad q(u)=\sum_i u(a_i)^2\quad(u\in V^*).
\end{equation}
The quadratic form $q$ has rank $p-2$. Moreover
$\dim L^{\circ2}=2p-1$; every preserving physical cubic weight is a scalar
multiple of the sheet-sign vector.
\end{proposition}
\begin{proof}
The sheets lie in the disjoint affine hyperplanes $\sigma=1,2$.
Within a sheet the points are distinct, and their differences span
$\ker\sigma$; $t$ supplies the remaining direction. Since $\sum_i a_i=0$,
translation preserves the first two moments of this $p$-point set.
The cubic expansion gives~\eqref{eq:translation-phase}.
Identify $V^*$ with $K=\{u\in\Fp^p:\sum_i u_i=0\}$.
Then $q(u)=\sum_i u_i^2$; its radical on $K$ is $\langle\one\rangle$,
so its rank is $p-2$ and the displayed cubic is nonzero.

For rigidity, let weights $\alpha_i,\beta_i$ on the two sheets annihilate
every affine quadratic. Put $\gamma=\alpha+\beta$ and $B=\sum_i\beta_i$.
The linear moment says $\gamma+Be_0\in\langle\one\rangle$.
The quadratic moment on $K$ has symmetric matrix
\[
 A=\operatorname{diag}(\gamma)+e_0\beta^T+\beta e_0^T+Be_0e_0^T.
\]
A symmetric form vanishing on $K\times K$ is
$\one r^T+r\one^T$ for some $r$. For distinct $i,j\ne0$, comparison gives
$r_i+r_j=0$. There are at least four such indices, whence all $r_i=0$
for $i\ne0$. Diagonal comparison gives $\gamma_i=0$ there; the linear
moment then forces $\gamma_0=-B$. Off-diagonal comparison gives
$\beta_i=r_0$ for $i\ne0$, and the $(0,0)$ entry gives $\beta_0=r_0$.
Thus $B=pr_0=0$, $\gamma=0$, and $\alpha=-\beta$ are constant weights.
The annihilator of $L^{\circ2}$ is therefore exactly the sign line.
\proves{prop:translation}\uses{thm:trade-code}
\end{proof}

\subsection{The conic configurations}
For the two symmetric examples take $Q=XZ-Y^2$ and label its points by
$a\mapsto(1:a:a^2)$ and $\infty\mapsto(0:0:1)$.
For a perfect matching $M$ of $\mathbb P^1(\Fp)$ define
\[
 P_M=\prod_{\{a,b\}\in M}L_{ab},\quad
 L_{ab}=abX-(a+b)Y+Z,\quad L_{a\infty}=aX-Y.
\]
Use base matchings
\[
 M_7=\{02,14,3\infty,56\},\qquad
 M_{11}=\{01,25,37,49,68,10\infty\}.
\]
Their $\PGL_2(p)$ orbits have $2p$ members. The determinant square class
of a transformation taking the base matching to $M$ specifies its sign.
The coefficient vector of $(P_M-P_{M_p})/Q$ specifies $x_M$.
The exact division, sign consistency, and moment conditions are part of the
finite input verification. This is the configuration studied in
\emph{Quadratic Trade Rigidity and Cubic Orientation in Conic Matching
Quotients}~\cite{tradegeometry}.

\begin{proposition}[Finite conic input]\label{prop:conic-input}
\coverage{absent}\uses{thm:trade-code}\evidence{conic-input}
For the stated matchings the construction has distinct points,
$\dim L=p$, and $\dim L^{\circ2}=2p-1$. It gives
$[[14,6,2]]_7$ and $[[22,10,2]]_{11}$.
In invertible logical coordinates the phases are
\begin{equation}\label{eq:normal-forms}
 F_7=4sI+3J,\qquad F_{11}=4sI_2+2I_3.
\end{equation}
Here $I=ae-4bd+3c^2$ and
$J=\det\left(\begin{smallmatrix}a&b&c\\b&c&d\\c&d&e\end{smallmatrix}\right)$
are binary-quartic invariants. For the binary octavic
$f(S,T)=\sum_{i=0}^8\binom8i z_iS^{8-i}T^i$,
$I_2=(f,f)_8$ and $I_3=((f,f)_4,f)_8$ use normalized transvectants.
\end{proposition}
\begin{proof}[Finite verification]
Appendix~\ref{sec:finite-data} specifies the matrices and substitutions.
Exact finite-field expansion checks the moments,
ranks and both polynomial identities. The signed-trade theorem then supplies
the code and gate. These are finite polynomial and linear-algebra checks;
the evidence record gives the independent reconstruction and its limits.
\proves{prop:conic-input}\uses{thm:trade-code}\evidence{conic-input}
\end{proof}
The invariant here is the logical phase itself. This differs from using
invariant theory to study weight enumerators~\cite{kalra}; it is not a claim
that invariant-defined gates in general originate with these examples.

\section{The Hessian spectrum and product exclusion}\label{sec:spectrum}
Clifford operations permute Pauli operators up to phase. The moduli of a state's
Pauli expectations therefore survive every Clifford change of frame and can
obstruct conversion to independent resources. Their fourth moment will give a
product-state obstruction; their rare largest nontrivial values will later
obstruct a short sum-of-cubes synthesis.

For a homogeneous cubic $F$ on $\Fp^k$, put
$|F\rangle=p^{-k/2}\sum_u\omega^{F(u)}|u\rangle$ and
$H_F(v)=(\partial_i\partial_jF)(v)$. We use natural logarithms and define
\[
 M_\alpha(|\psi\rangle)=\frac{1}{1-\alpha}
 \log\left(p^{-k}\sum_{v,w}|\langle\psi|X(v)Z(w)|\psi\rangle|^{2\alpha}\right)
 \quad(\alpha>0,\ \alpha\ne1).
\]
These stabilizer R\'enyi entropies quantify nonstabilizerness, usually called
magic, with the normalization in
\cite{leone,wangli}; they vanish on pure stabilizer states and are additive
under tensor products. The following quadratic-sum calculation makes them
accessible without enumerating all $p^{2k}$ Paulis.

\begin{theorem}[Rank formula]\label{thm:hessian}
\coverage{absent}
Let $p\geq5$, and set $r(v)=\rk H_F(v)$ and
$N_r=\#\{v:r(v)=r\}$. Then
\begin{equation}\label{eq:pauli-spectrum}
 |\langle F|X(v)Z(w)|F\rangle|=
 \begin{cases}p^{-r(v)/2},&w\in\im H_F(v),\\0,&\text{otherwise.}\end{cases}
\end{equation}
Consequently
\begin{equation}\label{eq:rank-entropy}
 M_\alpha(|F\rangle)=\frac{1}{1-\alpha}
 \log\left(p^{-k}\sum_{r=0}^kN_rp^{(1-\alpha)r}\right).
\end{equation}
\end{theorem}
\begin{proof}
The expectation is $p^{-k}\sum_u\omega^{F(u)-F(u+v)+w^Tu}$.
Writing $c_v=\nabla F(v)$, the exponent, up to a constant, is
$-\tfrac12u^TH_F(v)u+(w-c_v)^Tu$.
The sum over the radical vanishes unless $w-c_v\in\im H_F(v)$.
Euler's identity $H_F(v)v=2c_v$ makes this equivalent to the condition
in~\eqref{eq:pauli-spectrum}. On a complementary subspace diagonalize the
quadratic form. Each nonzero quadratic coefficient contributes a Gauss sum
of modulus $\sqrt p$: its squared modulus follows by replacing a pair
$(x,y)$ with $(x-y,x+y)$ in the double sum, an invertible substitution.
The radical contributes $p^{k-r(v)}$, proving the modulus formula.
There are $p^{r(v)}$ admissible $w$ for each $v$; summing their powers proves
\eqref{eq:rank-entropy}.
\proves{thm:hessian}
\end{proof}

Kagamihara--Tsuchiya~\cite[Theorem 1]{kagamihara} give the binary
rank-reduction formula, with $C(x)+C(x)^T$ replacing the odd-characteristic
Hessian. Chen--Yan--Zhou~\cite{chen} supply related hypergraph-state moment
methods. We use the established rank-reduction method in odd characteristic;
the contribution below is the geometric evaluation and its resource consequences.
For example, a single-qudit cubic $F(u)=u^3$ has rank zero at $u=0$ and rank
one at all other points. Thus $N_0=1$, $N_1=p-1$, and
$M_2=\log(p^2/(2p-1))$. Tensor products add these entropies. The same formula
now turns the multivariate rank counts into an exact resource comparison.

\begin{proposition}[Exact conic spectra]\label{prop:spectra}
\coverage{absent}\uses{prop:conic-input,thm:hessian}\evidence{rank-census}
The nonzero entries of the vector rank distributions for the two conic phases are
\[
\begin{array}{c|rrrrr}
p=7:\ r&0&3&4&5&6\\\hline
N_r&1&48&2940&26502&88158
\end{array}
\]
\[
\begin{array}{c|rrrrrrr}
p=11:\ r&0&5&6&7&8&9&10\\\hline
N_r&1&120&14520&80520&21442960&2387485210&23528401270.
\end{array}
\]
In each case the minimum positive rank is $(p-1)/2$, attained on exactly
$p+1$ projective points. In the binary-form coordinates these are $s=0$ and
the perfect $(p-3)$rd powers. In particular,
\[
 M_2(F_7)=\log\frac{1977326743}{78835},\qquad
 M_2(F_{11})=\log\frac{11^{19}}{7151076991}.
\]
\end{proposition}
\begin{proof}[Finite verification and deduction]
The census visits vectors with first nonzero coordinate one. Since
$H_F(av)=aH_F(v)$ for $a\ne0$, every projective representative contributes
$p-1$ vectors; the zero vector is added separately. The domains have
$19608$ and $2593742460$ representatives. The $p=7$ census has an independent
implementation; the $p=11$ exhaustive execution is checked against it in the
smaller case, not independently duplicated at its full size.
Substitution verifies the $p+1$ perfect-power points. The census has exactly
that many minimum-rank lines, so this containment is an equality over $\Fp$.
Equation~\eqref{eq:rank-entropy} gives the entropy values.
\proves{prop:spectra}\uses{thm:hessian}\evidence{rank-census}
\end{proof}

\begin{lemma}[Product ceiling]\label{lem:product-ceiling}
\coverage{absent}
For a pure state in dimension $D$ with an orthogonal unitary Pauli basis,
$M_2\leq\log((D+1)/2)$. Thus a pure product on blocks of $a$ and $k-a$
prime-dimensional qudits satisfies
\begin{equation}\label{eq:product-bound}
 M_2\leq\log\frac{(p^a+1)(p^{k-a}+1)}4.
\end{equation}
\end{lemma}
\begin{proof}
Put $x_P=|\langle P\rangle|^2$. Purity gives $\sum_Px_P=D$ and
$x_I=1$. Cauchy--Schwarz on the $D^2-1$ other entries yields
$\sum_Px_P^2\geq1+(D-1)^2/(D^2-1)=2D/(D+1)$.
Insert this into the definition of $M_2$. For a product the Pauli sum
factorizes, so the entropies add.
\proves{lem:product-ceiling}
\end{proof}
The product-ceiling method has predecessors in the characteristic-function
and qudit entropy literature~\cite{dai,wangli}. The bound need not be attained.

\begin{corollary}[Clifford-product exclusion]\label{cor:product-exclusion}
\coverage{absent}\uses{prop:spectra,lem:product-ceiling}\evidence{magic-comparison}
The ten-qudit state $|F_{11}\rangle$ is not Clifford-equivalent to a product
across any nontrivial bipartition. The same conclusion holds for the dihedral
and Borel six-qudit conic trades over $\mathbb F_7$ in
Appendix~\ref{sec:finite-data}. This entropy test alone does not establish
the conclusion for the six-qudit Clebsch state $|F_7\rangle$.
\end{corollary}
\begin{proof}
For fixed $k$, the largest right side of~\eqref{eq:product-bound} occurs at
$a=1$ or $k-1$: the variable part $p^a+p^{k-a}$ increases toward the
endpoints. At $p=11,k=10$ the bound is $22.6797\ldots$, below
$22.8695\ldots=M_2(F_{11})$. At $p=7,k=6$ it is $10.4228\ldots$;
the other two trades give $10.4246\ldots$ and $10.4840\ldots$, whereas
$M_2(F_7)=10.1299\ldots$. These comparisons are checked as exact rational
inequalities before taking logarithms. Clifford conjugation permutes the
Pauli moduli, so it preserves $M_2$.
\proves{cor:product-exclusion}\uses{prop:spectra,lem:product-ceiling}\evidence{magic-comparison}
\end{proof}

Two disjoint $xyz$ phases at $p=7$ have
$M_2=8.8047\ldots$, also below $F_7$.
The trace cubic $F(u)=\Tr_{\mathbb F_{p^k}/\Fp}(u^3)$ has full Hessian rank
at every nonzero $u$, by nondegeneracy of the trace pairing, and hence
$M_2=\log(D^2/(2D-1))$, $D=p^k$. These finite-field cubic fiducials belong
to the Alltop--MUB construction~\cite{alltop,klappenecker,damski}.
Feng--Luo~\cite[Proposition 2 and Appendix C]{feng} prove single-prime
L1 optimality among diagonal gates. Neither that result nor the L1 ceiling
of~\cite{damski} is an unrestricted $M_2$ maximality theorem.
The conjecture in~\cite{knipfer} concerns two prime-dimensional qudits;
we do not extend its attribution to arbitrary $k$.

\section{A bounded factory comparison}\label{sec:factory}
The resource-count question is different from the entropy question. We compare
preparing and purifying the whole six-qudit phase with purifying independent
single-qudit primitives and then compiling it. This is the synthillation
strategy of Campbell--Howard~\cite{campbellhoward}; Li--Yeh~\cite[Section 4.4
and Appendix C.1]{liyeh} already connects signed physical qutrit gates to a
coupled logical resource and its distillation. Our claim is the following
finite, same-target comparison.

\subsection{Noise and allowed operations}
Every raw input is
\begin{equation}\label{eq:noise}
 \rho_c=(1-\delta)|M_c\rangle\langle M_c|+
 \frac{\delta}{6}\sum_{e\ne0}Z(e)|M_c\rangle\langle M_c|Z(-e),
 \qquad |M_c\rangle=7^{-1/2}\sum_x\omega^{cx^3}|x\rangle.
\end{equation}
Inputs are independent; stabilizer preparation, Clifford operations,
measurements, feed-forward and storage are ideal. Uniformity of the six
nonzero labels is an assumption, not a consequence of diagonal twirling.
Every cost below counts raw magic inputs only; the ideal operations are not
priced, so the comparison is not a total implementation-overhead statement.

The native factory prepares $|+_L\rangle=7^{-7/2}\sum_{z\in L}|z\rangle$,
injects the fourteen signed cubic gates, measures the code stabilizers,
and decodes on acceptance. Injection is deterministic: controlled subtraction
from data $x$ to a magic ancilla, followed by ancilla outcome $m$, supplies
phase $c(x+m)^3$. A quadratic Clifford correction removes the $m$-dependent
terms. An ancilla error $Z(e)$ becomes $Z(e)$ on the data. No injection outcome
is discarded. For the resulting physical error vector $e$, acceptance is
$\one^Te=0$ and the logical error label is $E^Te$; the error is harmless
exactly when $e\in L^\perp$.

The separate architecture distills each primitive independently with any
finite concatenation of the three modules in Table~\ref{tab:modules}, then
uses deterministic linear-form cubic synthesis without final postselection.
Outputs may be stored separately and successful attempts retained. Module
choices and depths may differ between terms. A final Clifford is allowed if
its input is a homogeneous cubic-phase state on the same six qudits.
Ancilla-assisted general synthesis, pooled multi-output distillers and joint
synthillation are outside this menu.

\begin{table}[ht]
\centering\small
\begin{tabular}{lllll}
\toprule
Module & points & X space & logical row & cubic weights\\
\midrule
Weighted four-site & $0,1,2,3$ & $\langle1\rangle$ & $x$ & $(-1,3,-3,1)$\\
QRM six-site & $\mathbb F_7^*$ & $\langle x\rangle$ & $1$ & all $1$\\
RS seven-site & $\mathbb F_7$ & $\langle1,x\rangle$ & $x^2$ & all $1$\\
\bottomrule
\end{tabular}
\caption{The three allowed one-output distillers. In each row
$L=S+\langle\ell\rangle$ and the logical coefficient is $-1$;
coordinate negation supplies $+1$. Parameters are respectively
$[[4,1,2]]_7$, $[[6,1,2]]_7$, and $[[7,1,3]]_7$.}\label{tab:modules}
\end{table}
Here QRM means quantum Reed--Muller and RS means Reed--Solomon.
The six-site row is the prime-qudit QRM construction~\cite{cab,campbell2014};
the length-seven RS row is checked separately. We distinguish raw supplies.
In the \emph{sign-only menu}, only $c=\pm1$ has unit cost. The four-site
module costs eight raw injections because each $\pm3$ gate uses three signs.
Its two middle error rates are
$\delta_3=(6/7)[1-(1-7\delta/6)^3]$.
In the \emph{weighted relaxation}, every nonzero cubic coefficient has unit
raw cost and error $\delta$. This relaxation favors the separate architecture.

Define the weighted cubic rank $r(F)$ as the least number of nonzero terms in
$F=\sum_i c_i\ell_i^3$, with $c_i\in\Fp^*$ and nonzero linear forms $\ell_i$
over the same field. In linear-form synthesis each term consumes a primitive.
The spectral calculation therefore has a second use: it limits how few
independently prepared primitives can realize the target.

\begin{lemma}[Cubic synthesis lower bound]\label{lem:waring}
\coverage{absent}\uses{prop:spectra}\evidence{rank-census,conic-input}
The weighted linear-form cubic ranks obey
$9\leq r(F_7)\leq13$ and $15\leq r(F_{11})\leq21$.
The rank-nine lower bound also holds for any homogeneous cubic phase on six
qudits Clifford-equivalent to $|F_7\rangle$.
\end{lemma}
\begin{proof}
In a shortest decomposition $F=\sum_{i=1}^r c_i\ell_i^3$, the coefficient
vectors are projectively distinct, and they span the logical space: a nonzero
common zero $v$ of all $\ell_i$ would have Hessian rank zero, but the census
has $N_0=1$.
The Hessian at $v$ has rank at most the number of $\ell_i(v)\ne0$.
For $F_7$, killing five independent forms gives $r\geq8$.
If $r=8$, every six coefficient vectors must be independent, or a common
annihilator would have Hessian rank at most two. The $\binom85=56$
five-subsets then span distinct hyperplanes, producing 56 minimum-rank lines.
The census has only eight. Thus $r\geq9$.
For $F_{11}$ the same argument gives $r\geq14$ and excludes equality:
$\binom{14}9=2002>12$ minimum-rank lines would be needed.
The raw signed decompositions have one zero row, giving the upper bounds.
For a cubic phase the Pauli-modulus multiplicity at $p^{-r/2}$ is $N_rp^r$,
so the histogram recovers the rank census. Clifford invariance therefore
preserves every input to the lower-bound argument.
\proves{lem:waring}\uses{prop:spectra}\evidence{rank-census,conic-input}
\end{proof}
A weighted cube is implemented by Clifford computation of $\ell_i$, one
weighted cubic gate, and uncomputation. This signature-tensor synthesis
framework is explicit in Heyfron--Campbell~\cite[Section IV]{heyfron};
the rank bound is for that model, not arbitrary circuits.


The comparison uses three steps. Native error detection brings the block error
below the raw input error. Independence then prevents any unpurified synthesis
term from meeting that target, even if errors from different terms can cancel.
Finally, the rank bound forces at least nine independently purified terms.
The proof retains the exact interval bounds and both raw-resource supplies.

\begin{theorem}[Same-target resource comparison]\label{thm:factory}
\coverage{absent}\uses{thm:trade-code,lem:waring}\evidence{factory}
Under~\eqref{eq:noise}, for every $0<\delta\leq0.01$, let $q$ be the
native factory's output block infidelity. Its expected raw consumption is
less than $16.116$. Every separate construction in the sign-only menu
with block infidelity at most $q$ costs at least $54$ raw inputs per output;
in the weighted relaxation it costs at least $36$.
Feasible separate constructions meeting this target exist throughout the interval.
\end{theorem}
\begin{proof}
Write $a=1-\delta$ and $P$ for native acceptance. Single errors are rejected.
Every accepted weight-two error is harmful, since the evaluation points are
distinct. Their probability is $(91/6)\delta^2a^{12}$.
For a fixed error support of size at least two, fixing all but its last
nonzero label leaves at most one of six choices that accepts. A union bound
on pairs therefore gives
\begin{equation}\label{eq:native-bounds}
 a^{14}\leq P\leq1-14\delta a^{13},\qquad
 \frac{91\delta^2a^{12}}{6P}\leq q\leq\frac{91\delta^2}{6a^{14}}.
\end{equation}
Bernoulli's inequality implies
\[
 a^{12}(1+14\delta a)\geq(1-12\delta)(1+14\delta-14\delta^2)
 =1+\delta(2-182\delta+168\delta^2)\geq1.
\]
Hence $a^{12}\geq1-14\delta a^{13}\geq P$,
so $q\geq(91/6)\delta^2$, while
$q/\delta\leq(91/600)/(0.99)^{14}<1$.
Also $14/P\leq14/(0.99)^{14}<16.116$.

A noisy primitive on a nonzero form with coefficient vector $b$ contributes
an independent error label $eb$. For probability distributions on the
additive logical-label group,
$\|\mu*\nu\|_\infty\leq\|\mu\|_\infty$, since each convolution value
is an average. An unpurified raw term therefore forces final error at least
$\delta$, even with cancellations or a fixed Pauli correction.
The states $Z(\eta)|F\rangle$ are orthogonal, so this probability is exactly
a block-infidelity statement. An allowed final Clifford maps the ideal and
errored outputs together and preserves their overlaps. Since $q<\delta$, all synthesis terms must
be purified. Each independent tree has at least six raw leaves in the
sign-only menu and four in the weighted menu. Retrying cannot lower those
counts. Lemma~\ref{lem:waring} requires at least nine terms, giving $54$
and $36$.

For feasibility use the thirteen-term signed decomposition. A weighted
four-site primitive has error at most $\delta^2/a^4$, so thirteen give
block error at most $13\delta^2/(0.99)^4<(91/6)\delta^2$.
For the seven-site module harmful accepted errors have weight at least three,
giving error at most $35\delta^3/a^7$; thirteen give
$455\delta^3/a^7<(91/6)\delta^2$ throughout the interval.
A union bound suffices for these feasible upper estimates; the lower bound
already allowed cancellations.
\proves{thm:factory}\uses{thm:trade-code,lem:waring}\evidence{factory}
\end{proof}

At $\delta=0.01$ the exact enumerator calculation gives the comparison below.
The separate rows are feasible constructions, not claims to attain their
menu's optimum.
\begin{center}
\begin{tabular}{lrr}
\toprule
Construction & block infidelity & expected raw inputs\\
\midrule
Native fourteen-input & $0.00159853$ & $16.0894$\\
Weighted four-site, thirteen terms & $0.00133141$ & $54.1275$\\
Sign-only seven-site, thirteen terms & $0.0000130920$ & $97.6325$\\
\bottomrule
\end{tabular}
\end{center}
For an ordinary enumerator $A_w$ of $L^\perp$, native good-output probability
before conditioning is $\sum_w A_wa^{14-w}(\delta/6)^w$.
Acceptance also has the closed form
$P=[1+6(1-7\delta/6)^{14}]/7$. Thus the calculation includes harmless
accepted errors, rather than treating every nonzero physical error as failure.

\section{A chordal restriction of the logical cubic}\label{sec:shadow}
The logical cubic also provides an interface with the invariant pencil in
\emph{Chordal and Conference Cubics: Reconstruction and a Residual
$C_2$-Torsor}~\cite{shadowgeometry}. A restriction of a logical polynomial
is a resource on fewer variables; it does not by itself construct a subcode.

\begin{proposition}[Chordal shadow and sign choices]\label{prop:shadow}
\coverage{absent}\uses{prop:conic-input,thm:hessian}\evidence{shadow}
Over $\mathbb F_{11}$ the stabilizer $A_5$ of the base matching acts on the
logical space as $\one\oplus V_4\oplus V_5$. The restriction to its unique
five-dimensional constituent is, in the recorded coordinates,
\[
 8\det\begin{pmatrix}z_0&z_1&z_2\\z_1&z_2&z_3\\z_2&z_3&z_4\end{pmatrix}.
\]
It is a chordal member of the invariant pencil. At ranks $0,3,4,5$, the
chordal and conference censuses are respectively
\[
 (1,120,27720,133210),\qquad (1,300,22260,138490).
\]
Thus their five-qudit phase states are not Clifford-equivalent. At $p=7$ the restriction $s=0$ in
\eqref{eq:normal-forms} is $3J$.
Reversing the sheets sends $U_F$ to $U_F^{-1}$. The residual involution
exchanging the two chordal members acts as a linear Clifford change of frame
on the five shadow variables.
\end{proposition}
\begin{proof}[Finite identification and operational deduction]
Character idempotents in the reconstructed $A_5$ action have ranks $1,4,5$.
The recorded intertwiner identifies $V_5$ with the augmentation module on the
six matching pairs; polynomial substitution gives the Hankel determinant.
The two exact censuses then separate the resources by
Theorem~\ref{thm:hessian}. The $p=7$ restriction follows directly from
\eqref{eq:normal-forms}.
Sheet reversal changes every sign in~\eqref{eq:logical}, hence negates $F$.
For an invertible linear map $q$, the basis permutation $|z\rangle\mapsto|qz\rangle$
is Clifford: it sends $X(v)$ to $X(qv)$ and $Z(w)$ to $Z(q^{-T}w)$.
The recorded involution exchanges the two chordal polynomials, giving the stated
frame change. The identity is on the shadow, not a claim about all extensions
to the full ten-qudit gate.
\proves{prop:shadow}\uses{prop:conic-input,thm:hessian}\evidence{shadow}
\end{proof}
The singular scheme of the Hankel cubic is a rational normal quartic;
its twelve $\mathbb F_{11}$-points are not the entire scheme.
The existing lift tests rule out the identity-on-complement extension and
geometric $\PGL_2(11)$ lifts of the residual involution. They do not classify
arbitrary full-gate lifts. These boundaries are part of ledger N4.

\subsection*{Questions suggested by the construction}
The translation family shows that the parameters and cubic-weight rigidity
alone do not select the symmetric examples. Their distinction lies in the
phase geometry and its spectrum. It remains to explain the observed stopping
of the conic translation-class source beyond $p=13$, rather than merely
extend a matching enumeration. A second question is whether a balanced
$2p$-point trade can realize a trace-cubic resource on $p-1$ logical qudits.
Finally, the distance-three search inside the fixed $p=7$ evaluation space
fails, while the elementary distance-three RS examples do not establish the
same product exclusion. Finding a distance-at-least-three code with several
logical qudits and such an invariant phase is a separate construction problem.
The present factory comparison leaves pooled distillation and unrestricted
adaptive preparation open.

\appendix
\section{Finite data and verification boundary}\label{sec:finite-data}
The arguments proving the signed-moment dictionary, translation family,
Hessian formula, product ceiling and interval comparison are given in the
text. The conic matrices, polynomial identities, exact spectra and shadow
identification use finite computations. No Lean coverage is claimed for this
companion. The source annotations identify each dependence; the claim map
pins every theorem statement and records formal coverage as absent.

\subsection{Coordinates and exact inputs}
The signed evaluation matrices, sheet signs and raw cubics are fixed in
\texttt{supplement/reconstruction/common.py}; the augmentation embedding,
conference matrix, both chordal representatives and the involution of
Proposition~\ref{prop:shadow} are in
\texttt{verification/shadow-certificate.json}. For $p=7$ the coefficient order is
$X^2,XY,XZ,Y^2,YZ,Z^2$ and the logical substitution is
$u=(a,b,s+c,3s+c,d,e)$.
For $p=11$ the selected coefficient order is
\[
 X^4,X^3Y,X^3Z,X^2YZ,X^2Z^2,XY^2Z,XYZ^2,XZ^3,YZ^3,Z^4;
\]
the substitution is
$u=(z_0,z_1,7z_2,6z_3,3s+6z_4,s+3z_4,6z_5,7z_6,z_7,z_8)$.
The transvectant normalization is
\[
 (f,g)_r=\frac{(m-r)!(n-r)!}{m!n!}
 \sum_{j=0}^r(-1)^j\binom rj
 \frac{\partial^rf}{\partial S^{r-j}\partial T^j}
 \frac{\partial^rg}{\partial S^j\partial T^{r-j}},
\]
for binary forms of degrees $m,n$; all denominators used here are invertible.

For these matching configurations, the signed cubic line is also the Macaulay
inverse system of the Artinian reduction with Hilbert function
$(1,p-1,p-1,1)$~\cite[Corollary on self-association and cubic duality]{tradegeometry}.
Thus the same cubic enters the geometric duality and the linear-form
compilation problem.

For the other $p=7$ trades, one may reconstruct the positive and negative
sheets as translation classes of the following matchings. Entries are partner
lists on $0,\ldots,6,\infty$, with index $7$ representing infinity:
\[
\begin{array}{c|cc}
 &+&-\\\hline
\text{dihedral}&(1,0,5,4,3,2,7,6)&(1,0,6,4,3,7,2,5)\\
\text{Borel}&(1,0,4,6,2,7,3,5)&(1,0,6,5,7,3,2,4).
\end{array}
\]
Their vector censuses at ranks $0,3,4,5,6$ are respectively
\[
 (1,6,1092,20202,96348),\qquad (1,6,588,20286,96768).
\]
Their fourth-moment sums are $58711/16807$ and $55327/16807$.

\subsection{Classification and distance computations}
The conic search enumerates perfect matchings, identifies their translation
classes, and buckets classes by equal first and second moment tensors.
Pairs in a bucket are tested for a nonzero third-moment difference, then
quotiented by the stated affine action. The respective numbers of translation
classes for $p=5,7,11,13,17,19$ are
$3,15,945,10395,2027025,34459425$; the numbers of inequivalent hits under
$\operatorname{AGL}(1,p)$ are $0,3,2,1,0,0$.
This is a finite source classification, not a classification of all trades
or all Clifford-equivalent codes. Every hit has the signed moment condition,
$\dim L=p$ and Schur-square dimension $2p-1$.
The $p=13$ hit gives $[[26,12,2]]_{13}$; its stored 20000-vector Hessian
sample is not an exhaustive spectrum or a proved minimum-rank bound.
The new $p=11$ and $p=13$ X-distance estimates are only upper bounds.

The fixed-space distance search uses all $L'\subseteq L$ containing $\one$
with $\dim L'\geq3$. Put
$R(L')=\{s\in L':T(s,L',L')=0\}$.
It suffices to test the maximal admissible X-space $R(L')$.
A weight-two $e_i-e_j$ is undetected exactly when columns $i,j$ of this
space agree; it is harmless exactly when those columns of $L'$ agree.
Thus distance at least three requires equality of the two column partitions.
Reduced row-echelon enumeration of subspaces of $L/\langle\one\rangle$
visits $61927311$ candidates of dimensions two through six. None with a
nonzero restricted cubic passes the partition test. The one-dimensional
case also fails: a nonzero one-variable cubic has zero trilinear radical
on that quotient, so its only available X-space is $\langle\one\rangle$,
which does not separate its distinct evaluation columns.
This excludes the specified nonzero-cubic subcodes inside the fixed $L$;
it does not exclude other evaluation spaces, concatenation or code switching.
For comparison, evaluation on all of $\Fp$ with X-space
$\operatorname{RS}_2$ gives the length-$p$ distance-three rows
$[[7,1,3]]_7$, $[[11,3,3]]_{11}$ and $[[13,4,3]]_{13}$ in the recorded
polynomial-code search. They are distinct from Campbell's length-$(p-1)$
construction~\cite{campbell2014}; their checks belong to the strengthening
bundle.
The exhaustive negative has an internal count check and recorded tally,
but no independent full second implementation.

\subsection{Artifacts and reproducibility}
The machine-readable evidence registry in \texttt{verification/evidence.json}
records the exact repository paths, checksums, byte counts, replay commands,
search domains and limits. The inputs are the committed reconstruction,
spectrum, strengthening and shadow bundles, together with the factory benchmark. The factory certificate contains all native and
primitive enumerators, exact rational sample values and interval constants.
Its checks compare primal enumeration with MacWilliams transforms and
synthesis-kernel enumeration with the Fourier expression; no random seed or
floating-point decision enters that benchmark.

The lightweight draft check verifies statement identities, annotation links,
input hashes, census arithmetic and the exact product and interval inequalities.
A fresh standard-library check also re-expands both coordinate normal forms,
checks both conic moment and Schur-square identities, reproduces the complete
$p=7$ census, and checks seven primes in the translation family.
A separate check reproduces the actual shadow normalization and compares all
$1320$ geometric actions with the residual involution.
It does not rerun the multi-billion-point $p=11$ census or the matching
exhaustions. Those remain explicitly attributed to the recorded executions;
full commands and their runtime estimates accompany the evidence registry.
The accompanying supplement contains the inputs and programs needed for these
checks. The lightweight check does not independently reproduce the large searches.

\begin{thebibliography}{99}

\bibitem{watson} Qudit Colour Codes and Gauge Colour Codes in All Spatial Dimensions. \href{https://arxiv.org/abs/1503.08800}{arXiv:1503.08800}.

\bibitem{haah} Towers of generalized divisible quantum codes. \href{https://arxiv.org/abs/1709.08658}{arXiv:1709.08658}.

\bibitem{cab} Magic state distillation in all prime dimensions using quantum Reed-Muller codes. \href{https://arxiv.org/abs/1205.3104}{arXiv:1205.3104}.

\bibitem{krishna} Towards low overhead magic state distillation. \href{https://arxiv.org/abs/1811.08461}{arXiv:1811.08461}.

\bibitem{kalra} Invariant Theory, Magic State Distillation, and Bounds on Classical Codes. \href{https://arxiv.org/abs/2501.10163}{arXiv:2501.10163}.

\bibitem{leone} Stabilizer R\'enyi Entropy. \href{https://arxiv.org/abs/2106.12587}{arXiv:2106.12587}.

\bibitem{wangli} Stabilizer R\'enyi entropy on qudits. \href{https://doi.org/10.1007/s11128-023-04186-9}{10.1007/s11128-023-04186-9}.

\bibitem{kagamihara} Stabilizer R\'enyi entropy of 3-uniform hypergraph states. \href{https://arxiv.org/abs/2602.23687}{arXiv:2602.23687}.

\bibitem{chen} Magic of quantum hypergraph states. \href{https://arxiv.org/abs/2308.01886}{arXiv:2308.01886}.

\bibitem{dai} Detecting Magic States via Characteristic Functions. \href{https://doi.org/10.1007/s10773-022-05027-8}{10.1007/s10773-022-05027-8}.

\bibitem{alltop} Complex sequences with low periodic correlations. \href{https://doi.org/10.1109/TIT.1980.1056185}{10.1109/TIT.1980.1056185}.

\bibitem{klappenecker} Constructions of Mutually Unbiased Bases. \href{https://arxiv.org/abs/quant-ph/0309120}{arXiv:quant-ph/0309120}.

\bibitem{damski} Product Weyl-Heisenberg covariant MUBs and Maximizers of Magick. \href{https://arxiv.org/abs/2603.15550}{arXiv:2603.15550}.

\bibitem{feng} Optimality of the Howard-Vala T-gate in stabilizer quantum computation. \href{https://doi.org/10.1088/1402-4896/ad80e7}{10.1088/1402-4896/ad80e7}.

\bibitem{knipfer} Analytical Landscape of Maximal Magic for Two-Qutrit States and Beyond. \href{https://arxiv.org/abs/2607.07197}{arXiv:2607.07197}.

\bibitem{campbellhoward} A unified framework for magic state distillation and multi-qubit gate synthesis with reduced resource cost. \href{https://arxiv.org/abs/1606.01904}{arXiv:1606.01904}.

\bibitem{liyeh} Transversal AND in Quantum Codes. \href{https://arxiv.org/abs/2603.04548}{arXiv:2603.04548}.

\bibitem{heyfron} A quantum compiler for qudits of prime dimension greater than 3. \href{https://arxiv.org/abs/1902.05634}{arXiv:1902.05634}.

\bibitem{campbell2014} Enhanced fault-tolerant quantum computing in d-level systems. \href{https://arxiv.org/abs/1406.3055}{arXiv:1406.3055}.

\bibitem{tradegeometry} Quadratic Trade Rigidity and Cubic Orientation in Conic Matching Quotients. Companion manuscript, repository source \texttt{papers/clebsch-factorization/}.

\bibitem{shadowgeometry} Chordal and Conference Cubics: Reconstruction and a Residual $C_2$-Torsor. Companion manuscript, repository source \texttt{papers/chordal-conference-reconstruction/}.

\bibitem{mds} Diagonal Isoduality and Transversal Clifford Groups of MDS--CSS Codes.
Companion manuscript, repository source \texttt{papers/mds\_css\_transversal\_groups/}.

\end{thebibliography}

\end{document}
